% !TeX root = ../../Book5.tex
% 纲要标题：### 5.4 初等子群检测与 Brauer 整诱导的证明
% 纲要位置：计划纲要.md，第 3223 行。

\section{初等子群检测与 Brauer 整诱导的证明}
\label{b5:sec:0504}

% TODO: 按以下纲要要点撰写本节。

% 纲要内容（仅作写作计划，尚非教材正文）：
%
% * 把诱导像组织为复表示环中的加法子群，利用限制、诱导与张量恒等式控制其环论行为
% * 建立初等子群上的限制检测与整性引理，分素数处理局部化后的整数系数问题，再由有限生成阿贝尔群的局部检测回到整数系数
% * 分别完成“存在某个整数倍的诱导表达”与“无需整体分母的整诱导表达”，指出第二步是超出 Artin 有理诱导的实质
% * 解释初等群上的表示如何进一步化为从适当子群一维表示的诱导；逐项说明所用子群仍处于允许的初等范围
% * 用小群的虚特征标表达说明整数系数可以为负，不把虚表示恒等式误解为实际表示的直接和分解
%

\subsection{诱导像与复表示环}
\label{b5:subsec:0504:01}

待撰写。

\subsection{限制检测、局部化与整性引理}
\label{b5:subsec:0504:02}

待撰写。

\subsection{从有理诱导到整系数诱导}
\label{b5:subsec:0504:03}

待撰写。

\subsection{初等群表示与一维表示的诱导}
\label{b5:subsec:0504:04}

待撰写。

\subsection{负系数与虚特征标的例子}
\label{b5:subsec:0504:05}

待撰写。
